\section{Finite trace recovery}\label{sec:trace}

\subsection{The source contract}

Write \(T=T^+\), \(T^*=T^-\), and \(e\) for the compact averaging operator.
The source predicate \code{LineRelation d c A Tp Tm e} provides
\begin{equation}\label{eq:line}
 d>0,\qquad c>0,\qquad e^2=e,\qquad
 TT^*=d(\Id+ce),\qquad AT=TA,\qquad AT^*=T^*A.
\end{equation}
Its existing trace-recovery theorem compares two such systems from
\[
 \tr(A^{l+1}T^{j+1})=\tr(A'^{\,l+1}T'^{\,j+1})
 \qquad(l,j\in\N).
\]
It passes to generalized eigenspaces of \(A\), determines characteristic polynomials
of the restricted raising operators, reads off determinant norms, recovers ranks of
the averaging projector, and reassembles the trace.

The following theorem needs a strictly smaller collection of moments.

\begin{theorem}[Finite positive moments suffice]\label{thm:finite-trace}
Let \(V,V'\) be finite-dimensional complex inner-product spaces. Let
\((A,T,e)\) and \((A',T',e')\) satisfy \eqref{eq:line} with the same positive
numbers \(d,c\). Let
\[
 S=\spec(T)\cup\spec(T'),\qquad s=\#S,
\]
where eigenvalues are counted distinctly in \(S\). If
\begin{equation}\label{eq:finite-moments}
 \tr(AT^j)=\tr(A'T'^{\,j})\qquad(1\leq j\leq s),
\end{equation}
then
\[
 \tr(Ae)=\tr(A'e').
\]
Both raising operators are invertible, so \(0\notin S\), and
\(s\leq\dim V+\dim V'\). Neither raising operator is required to be normal
or diagonalizable.
\end{theorem}

The caller need not compute \(S\). Providing the moments through
\(\dim V+\dim V'\), or providing all positive moments with only one index,
is sufficient. Thus an implementation can initially retain a convenient
unbounded positive-power premise while removing all higher powers of \(A\).

\subsection{A polynomial that represents the projector inside traces}

For \(z\ne0\), set
\begin{equation}\label{eq:log-weight}
 f(z)=\frac{\log(\abs{z}^2/d)}{\log(1+c)}
     =\frac{2\log\abs z-\log d}{\log(1+c)}.
\end{equation}
The denominator is nonzero because \(c>0\). Interpolate \(f\) on \(S\), together
with the value \(0\) at the point \(0\). The resulting polynomial \(p\) satisfies
\begin{equation}\label{eq:polynomial}
 p(0)=0,\qquad \deg p\leq s,\qquad p(\lambda)=f(\lambda)\quad(\lambda\in S).
\end{equation}
Equivalently, for nonempty \(S\), one can use the explicit formula
\[
 p(X)=X\sum_{\lambda\in S}\frac{f(\lambda)}{\lambda}
          \prod_{\substack{\mu\in S\\\mu\ne\lambda}}
          \frac{X-\mu}{\lambda-\mu}.
\]
Only distinct nodes are interpolated. No Hermite interpolation of derivatives
is needed, even when \(T\) has nontrivial Jordan blocks.

\begin{lemma}[Spectral trace identity]\label{lem:spectral-certificate}
Under \eqref{eq:line}, if \(p(\lambda)=f(\lambda)\) for every eigenvalue of \(T\),
then
\begin{equation}\label{eq:certificate}
 \tr(Ae)=\tr\bigl(Ap(T)\bigr).
\end{equation}
\end{lemma}

\begin{proof}
First, \(e=(d^{-1}TT^*-\Id)/c\) is self-adjoint, hence an orthogonal projection.
The identity \eqref{eq:line} shows that \(TT^*\) is positive definite, so \(T\)
is invertible. In dimension zero these statements and all subsequent trace
equalities have their usual empty-space interpretations.

Suppose first that \(A=A^*\). Decompose \(V\) orthogonally into its eigenspaces
\(H_\alpha\). Both \(T\) and \(T^*\) preserve every \(H_\alpha\), because both commute
with \(A\). Thus each \(H_\alpha\) reduces \(T\); the adjoint of the restriction of
\(T\) is the restriction of \(T^*\). The projection \(e\) also preserves the block.
Put
\[
 n_\alpha=\dim H_\alpha,\qquad
 r_\alpha=\rank(e|_{H_\alpha}),\qquad T_\alpha=T|_{H_\alpha}.
\]
Taking determinants of the restricted relation gives
\begin{equation}\label{eq:block-det}
 \abs{\det T_\alpha}^2
       =d^{n_\alpha}(1+c)^{r_\alpha}.
\end{equation}
For the eigenvalues of \(T_\alpha\), now counted with algebraic multiplicity,
\begin{align*}
 \sum_\lambda f(\lambda)
 &=\frac{2\sum_\lambda\log\abs\lambda-n_\alpha\log d}{\log(1+c)}\\
 &=\frac{\log\abs{\det T_\alpha}^2-n_\alpha\log d}{\log(1+c)}
 =r_\alpha.
\end{align*}
Every eigenvalue of a restriction to an invariant subspace belongs to the full
spectrum of \(T\). Triangularizing \(T_\alpha\) gives
\(\tr p(T_\alpha)=\sum_\lambda p(\lambda)=r_\alpha\).
This uses algebraic multiplicities, not diagonalizability. Consequently,
\[
 \tr(Ap(T))=\sum_\alpha\alpha\,\tr p(T_\alpha)
            =\sum_\alpha\alpha r_\alpha
            =\tr(Ae).
\]

For general \(A\), write
\[
 A=H+iK,\qquad H=\frac{A+A^*}{2},\qquad K=\frac{A-A^*}{2i}.
\]
The two commutation hypotheses for \(A\), after taking adjoints, imply that
\(A^*\) also commutes with \(T\) and \(T^*\). Hence \(H,K\) are self-adjoint
and commute with both. Apply the preceding argument twice, using the same
polynomial \(p\), and combine by complex linearity of trace.
\end{proof}

\begin{proof}[Proof of Theorem~\ref{thm:finite-trace}]
Choose \(p\) simultaneously on the union of spectra as in \eqref{eq:polynomial}.
Write \(p(X)=\sum_{j=1}^{s}b_jX^j\). Then
\[
 \tr(Ae)=\sum_{j=1}^{s}b_j\tr(AT^j)
        =\sum_{j=1}^{s}b_j\tr(A'T'^{\,j})
        =\tr(A'e').
\]
If \(S\) is empty, both spaces have dimension zero and the assertion is immediate.
\end{proof}

\subsection{Why the hypotheses and the trace formulation matter}

The assertion is generally false as an equality of operators \(e=p(T)\).
It is an equality of their trace pairings with the common commutant of
\(T\) and \(T^*\). This distinction also prevents the invalid inference that
\eqref{eq:line} makes \(T\) normal.

Here is an exact example. Set
\[
 T=\begin{pmatrix}1&3/2\\0&1\end{pmatrix},\qquad
 e=\begin{pmatrix}4/5&2/5\\2/5&1/5\end{pmatrix},
 \qquad d=\tfrac14,\quad c=15.
\]
Then \(e^2=e\) and \(TT^*=\tfrac14(\Id+15e)\).
The only eigenvalue of \(T\) is \(1\), so \(f(1)=1/2\), and
\(p(X)=X/2\) is a valid zero-constant interpolant. But \(e\ne p(T)\).
If
\[
 A=\begin{pmatrix}0&1\\0&0\end{pmatrix},
\]
then \(AT=TA\), while \(AT^*\ne T^*A\), and
\[
 \tr(Ae)=\tfrac25,\qquad \tr(Ap(T))=0.
\]
Thus the second commutation hypothesis cannot simply be discarded. The source
\code{LineRelation} already supplies it.

Nor does a single positive moment always suffice. With
\[
 T=\operatorname{diag}(1,2),\quad e=\operatorname{diag}(0,1),\quad
 A=\operatorname{diag}(-2,1),\quad d=1,\quad c=3,
\]
one has \(\tr(AT)=0\) but \(\tr(Ae)=1\).
The theorem retains enough moments for the number of distinct spectral nodes.
The archive checks both examples exactly over \(\mathbb Q\).

\subsection{The source-level gain}

In \code{LevelCut}, keep the local operators and the same line relation.
Replace the trace premise involving
\[
 \texttt{lineHecke}^{\,m+1}\cdot\texttt{offPlace}(a^{\,l+1})
\]
by one involving
\[
 \texttt{lineHecke}^{\,m+1}\cdot\texttt{offPlace}(a).
\]
The existing action identity specializes at \(l=0\).
In \code{CuspidalComparison}, the original cuspidal test \(a\) remains present;
the positive power of the auxiliary line operator still imposes the simple-zero
condition on the support. The outer theorem
\code{trace_eq_of_isCuspidalAt} can keep its existing signature.

This replacement does not ask the global group algebra to be closed under adjoints.
The decomposition \(A=H+iK\) occurs entirely inside the finite-dimensional
linear-algebra lemma. That avoids an unnecessary new assumption on the
adelic test-function algebra.
